% =============================================================================
% Temperature Control Case Study — Full Expanded Version
% IEEE-standard LaTeX, ready for Overleaf compilation
%
% Required packages (add to main thesis preamble if not already present):
%   \usepackage{algorithm}
%   \usepackage{algorithmic}
%   \usepackage{booktabs}
%   \usepackage{multirow}
%   \usepackage{graphicx}
%   \usepackage{amsmath,amssymb}
%   \usepackage{xcolor}
%   \usepackage{siunitx}   % for \num{}
% =============================================================================

\chapter{Temperature Control Case Study}
\label{ch:temperature-case-study}

This chapter presents a comprehensive experimental study of safe reinforcement
learning (RL) applied to multi-building HVAC temperature control in the
CityLearn simulation environment.  The task is framed as a Constrained Markov
Decision Process (CMDP) in which a centralized agent must minimise electricity
cost, grid stress, and carbon emissions while keeping indoor temperatures within
the occupant comfort band $[20\,^\circ\text{C},\, 26\,^\circ\text{C}]$ at all
times.

The chapter makes three main contributions.  \textbf{First}, it introduces the
SafeCityLearn framework---a modular simulation and training infrastructure that
couples CityLearn~v2~\cite{nweye2023citylearn} with the OmniSafe
constrained-RL backend~\cite{omnisafe2023}---and describes in detail the
CMDP formulation, observation design, and STEMS reward decomposition adopted
throughout.  \textbf{Second}, it provides an in-depth derivation and training
analysis of the Constrained Soft Actor-Critic with a smoothed Log-Barrier
penalty (\textbf{CSAC-LB})~\cite{zhang2024csaclb}, explaining every
algorithmic choice at the gradient level.  \textbf{Third}, it presents a
multi-algorithm, multi-seed benchmark comparing CSAC-LB against representative
safe-RL baselines---including Constrained Policy Optimisation
(\textbf{CPO})~\cite{achiam2017cpo}, First-Order Constrained Optimisation in
Policy Space (\textbf{FOCOPS})~\cite{zhang2020focops}, Constrained Update
Projection (\textbf{CUP})~\cite{yang2022cup}, and SAC-Lagrangian
(\textbf{SAC-Lag})~\cite{ha2020learning}---together with unconstrained and
rule-based reference policies.  Statistical robustness is quantified via
\num{3}--\num{6} random seeds per algorithm, and robustness under extreme
conditions is characterised via a 72-hour simulated heat-wave stress test.


% =============================================================================
\section{SafeCityLearn System Architecture}
\label{sec:temp-architecture}
% =============================================================================

All experiments are implemented on top of the \emph{SafeCityLearn} framework,
whose system architecture is depicted in Figure~\ref{fig:safecitylearn-arch}.
The framework forms a three-tier stack connecting a realistic building-physics
simulator to a state-of-the-art constrained-RL library through a modular
interfacing layer.

\begin{figure}[htbp]
  \centering
  \includegraphics[width=0.92\textwidth]{figures/Pasted_image.png}
  \caption{SafeCityLearn system architecture.
    \textbf{Top tier:} The Gymnasium~\cite{towers2023gymnasium} API provides
    the environment interface.  Two concrete environment branches are derived:
    (left) \texttt{CityLearn V2-V2G} for vehicle-to-grid battery control and
    (right) \texttt{CityLearn HVAC} for the temperature case study presented
    in this chapter.
    \textbf{Middle tier:} SafeCityLearn acts as the central integration hub.
    It registers each CityLearn variant as a named Gymnasium environment,
    applies reward shaping (STEMS decomposition), action masking, and
    observation normalisation, and exposes cost signals required by constrained
    RL.  The OmniSafe~\cite{omnisafe2023} CMDP backend provides all algorithm
    implementations through a unified trainer API, enabling fair algorithm
    comparisons with shared environment wrappers and random seeds.
    \textbf{Left branch (V2G task):} On-policy algorithms---PPO-Lag, CPO,
    FOCOPS, and Multi-Lagrangian PPO---are used for the vehicle-to-grid
    task, leveraging their lower memory footprint on the complex
    multi-constraint V2G CMDP.
    \textbf{Right branch (Temperature task):} CSAC-LB is the primary
    algorithm for the temperature case study due to its off-policy
    sample efficiency and conservative safety critic.
    \textbf{Bottom tier:} A three-level evaluation system produces
    cooling-season deterministic rollouts, multi-seed statistical summaries, and
    extreme-weather stress tests.}
  \label{fig:safecitylearn-arch}
\end{figure}

\paragraph{Design rationale.}
The separation between the environment tier and the algorithm tier is
intentional.  By routing all algorithms through the same SafeCityLearn wrapper,
the experimental results presented in Section~\ref{sec:temp-results} are
designed to differ primarily in the RL algorithm rather than in preprocessing,
reward scaling, or data normalisation.  The OmniSafe backend was selected because it
provides reference implementations for the main constrained-RL families used in
this thesis---including CSAC-LB, CPO, FOCOPS, CUP, SAC-Lag, PPO-Lag, and
unconstrained PPO---using a common neural network architecture, optimizer, and
hyperparameter convention.


% =============================================================================
\section{Methodology}
\label{sec:temp-methodology}
% =============================================================================

% -----------------------------------------------------------------------------
\subsection{Simulation Environment}
\label{sec:temp-env}

Experiments are conducted on CityLearn~v2~\cite{nweye2023citylearn}, a
multi-building energy management simulation platform that provides realistic
building energy dynamics and grid interactions.  The experimental dataset is
derived from real building energy consumption data in Travis County, Texas
(ASHRAE climate zone~2A: hot-humid).  For this case study, we use the
cooling-season subset corresponding to June--August, yielding a
\num{2207}-step hourly episode focused on summer HVAC operation rather than
year-round control.  This distinction matters for interpretation: all energy,
comfort, and cost figures reported in this chapter refer to the summer
cooling season only.

The case-study schema binds three residential buildings with LSTM-enabled
indoor temperature dynamics.  Each building is equipped with a cooling device,
electrical storage, and domestic hot water (DHW) storage.  A centralised
single-agent controller issues actions for all three buildings at hourly
resolution, yielding $T = 2{,}207$ hourly steps per episode.  The first
13 steps form an LSTM warm-up period during which comfort costs are disabled,
leaving 2{,}194 active evaluation steps.
The LSTM thermal model is a recurrent network trained on historical indoor
temperature data and conditioned on outdoor temperature, solar irradiance, and
occupancy; the 13-step warm-up allows it to initialise its hidden state from
the start-of-episode boundary conditions before the agent's actions begin
affecting indoor temperatures.

% -----------------------------------------------------------------------------
\subsection{CMDP Formulation}
\label{sec:temp-formulation}

The temperature control problem is formulated as a Constrained Markov Decision
Process (CMDP) $\mathcal{M} = (\mathcal{S}, \mathcal{A}, P, \mathcal{R},
\mathcal{C}, \gamma)$.

\paragraph{State space.}
The observation vector $\mathbf{s}_t \in \mathcal{S} \subset \mathbb{R}^{54}$
is constructed from normalised per-building, environmental, and temporal
features (Table~\ref{tab:observation-space}).  All features are normalised to
$[0, 1]$ by CityLearn's internal observation wrapper using dataset-derived
bounds.

\begin{table}[htbp]
\centering
\caption{Observation space composition ($\dim(\mathcal{S}) = 54$).  Each
  per-building feature is replicated for all $N=3$ buildings.}
\label{tab:observation-space}
\begin{tabular}{llc}
\toprule
\textbf{Category} & \textbf{Feature} & \textbf{Dim} \\
\midrule
\multirow{5}{*}{Building state}
  & Indoor dry-bulb temperature $T_t^{(i)}$ & $3$ \\
  & Temperature setpoint / reference & $3$ \\
  & Cooling device status & $3$ \\
  & Electrical storage state of charge & $3$ \\
  & DHW storage state of charge & $3$ \\
\midrule
\multirow{4}{*}{Environmental}
  & Outdoor dry-bulb temperature & $1$ \\
  & Direct solar irradiance & $1$ \\
  & Diffuse solar irradiance & $1$ \\
  & Solar generation $p_t^{(i)}$ & $3$ \\
  & Electricity price $v_t$ & $1$ \\
\midrule
\multirow{3}{*}{Temporal}
  & Hour of day (sine/cosine encoded) & $2$ \\
  & Day of year (sine/cosine encoded) & $2$ \\
  & Occupancy count per building & $3$ \\
\midrule
\multirow{2}{*}{Grid / district}
  & Net district electricity consumption & $1$ \\
  & District import power & $1$ \\
  & Net building consumption $e_t^{(i)}$ & $3$ \\
\midrule
\multicolumn{2}{l}{\textit{Remaining CityLearn default features}} & $\sim$20 \\
\bottomrule
\multicolumn{2}{l}{\textbf{Total}} & \textbf{54} \\
\end{tabular}
\end{table}

The state provides the agent with the information necessary to reason about
thermal dynamics (indoor and outdoor temperatures), energy markets (electricity
price), renewable availability (solar generation), and temporal patterns
(time encoding).  The sine/cosine temporal encoding enables the policy to learn
periodic pricing and occupancy patterns without requiring the network to
extrapolate across a discrete day-of-week feature.

\paragraph{Action space.}
In the \emph{cooling-only} interface adopted for this case study, the policy
controls a single continuous cooling action per building,
$\mathbf{a}_t \in [0,1]^3$, where $a_t^{(i)} = 0$ means no cooling and
$a_t^{(i)} = 1$ means maximum cooling for building~$i$.  Battery and DHW
storage actions are fixed to zero by the environment wrapper to isolate the
temperature control problem from storage management (Table~\ref{tab:action-space}).

\begin{table}[htbp]
\centering
\caption{Action space definition.  Only cooling actions are controlled by the
  RL policy; all other actuators are frozen to isolate the temperature-control
  problem.}
\label{tab:action-space}
\begin{tabular}{lccc}
\toprule
\textbf{Actuator} & \textbf{Range} & \textbf{Controlled by} & \textbf{Dim} \\
\midrule
Cooling device & $[0, 1]$ & RL policy & $3$ \\
Electrical storage & $0$ (fixed) & Wrapper & $3$ \\
DHW storage & $0$ (fixed) & Wrapper & $3$ \\
\midrule
\multicolumn{3}{l}{\textbf{Effective policy action dimension}} & \textbf{3} \\
\bottomrule
\end{tabular}
\end{table}

\paragraph{Reward function.}
Following the STEMS framework~\cite{zhang2025stems}, the reward is decomposed
into four components:
\begin{equation}
  R_t = R_t^{\text{economic}} + R_t^{\text{stability}} +
        R_t^{\text{renewable}} + R_t^{\text{comfort}}
  \label{eq:reward-decomposition}
\end{equation}

The \emph{economic} component minimises electricity import cost:
\begin{equation}
  R_t^{\text{economic}} = -\mu \cdot v_t \cdot \sum_{i=1}^{N} e_t^{(i)},
  \quad \mu = 1.0
  \label{eq:reward-economic}
\end{equation}
where $v_t$ is the electricity price and $e_t^{(i)}$ is the net electricity
consumption of building~$i$.  This term scales linearly with both price and
consumption, incentivising the agent to shift loads toward low-price periods.

The \emph{stability} component rewards smooth grid operation:
\begin{equation}
  R_t^{\text{stability}} = \underbrace{N \cdot \alpha_{\text{grid}} \left(
    1 - \left(\frac{\sum_i \max(0, e_t^{(i)})}{P_{\text{grid,max}}}\right)^2
  \right)}_{\text{grid load factor}}
  + \underbrace{\alpha_{\text{build}} \sum_{i=1}^{N} \left(
    1 - \frac{|e_t^{(i)}|}{P_{\text{build,max}}}
  \right)}_{\text{building load factor}}
  - \underbrace{\beta_{\text{ramp}} \sum_{i=1}^{N} \frac{|e_t^{(i)} - e_{t-1}^{(i)}|}{P_{\text{build,max}}}}_{\text{ramping penalty}}
  \label{eq:reward-stability}
\end{equation}
with $\alpha_{\text{grid}} = 0.5$, $\alpha_{\text{build}} = 0.3$,
$\beta_{\text{ramp}} = 0.2$.  The quadratic grid load factor penalises
high-demand peaks more strongly than low demand, matching the cost structure
of real-world electricity grids where peak demand incurs non-linear tariffs.

The \emph{renewable energy} component maximises self-consumption of solar
generation:
\begin{equation}
  R_t^{\text{renewable}} = \xi \sum_{i=1}^{N}
    \min\!\left(\frac{p_t^{(i)}}{p_t^{(i)} + \max(0, e_t^{(i)})},\; 1\right),
  \quad \xi = 0.6
  \label{eq:reward-renewable}
\end{equation}
where $p_t^{(i)}$ is the solar generation.  The ratio inside the $\min$ is the
solar self-sufficiency fraction: it approaches 1 when solar generation exceeds
net consumption and 0 when the building imports heavily from the grid.

The \emph{comfort} component penalises temperature deviations with
occupancy-aware weighting:
\begin{equation}
  R_t^{\text{comfort}} = -\lambda_{\text{indoor}} \sum_{i=1}^{N}
    w_i \cdot \Delta T_t^{(i)}
  \label{eq:reward-comfort}
\end{equation}
where $\Delta T_t^{(i)} = \max(T_t^{(i)} - T_{\max}, 0) + \max(T_{\min} -
T_t^{(i)}, 0)$ is the one-sided comfort deviation, $w_i = 1.0$ during
occupancy and $w_i = \eta = 0.1$ when vacant, and
$\lambda_{\text{indoor}} = 1.2$.  This reward gradient is \emph{complementary}
to---not a substitute for---the explicit CMDP cost constraint: it provides a
dense shaped signal that guides the policy toward the comfort band between
binary cost firings, analogous to reward shaping in potential-based
formulations~\cite{ng1999policy}.

\paragraph{Cost function.}
The CMDP cost signal is a district-level comfort violation indicator:
\begin{equation}
  C_t = w_{\text{comfort}} \cdot \sum_{i=1}^{N} \mathbf{1}\!\left[
    T_t^{(i)} < T_{\min} \;\text{or}\; T_t^{(i)} > T_{\max}
  \right], \quad w_{\text{comfort}} = 0.05
  \label{eq:cmdp-cost}
\end{equation}
During the 13-step LSTM warm-up, $C_t = 0$.  The constraint budget for the
CMDP optimisation problem is:
\begin{equation}
  J_C(\pi) = \mathbb{E}_\pi\!\left[\sum_{t=0}^{T} \gamma^t C_t\right] \leq d
  \label{eq:cmdp-budget}
\end{equation}
where $d$ is set per-algorithm (e.g.\ $d = 0.5$ for CSAC-LB, $d = 40.0$ for
CPO) to match the cost scale of each algorithm's internal constraint
representation.  The \emph{violation rate} metric reported throughout this
chapter is the fraction of active steps at which at least one building violates
the comfort band:
\begin{equation}
  V_{\text{rate}} = \frac{1}{T_{\text{active}}}
    \sum_{t=1}^{T_{\text{active}}} \mathbf{1}\!\left[
      \exists\, i : T_t^{(i)} < T_{\min}
      \;\text{or}\; T_t^{(i)} > T_{\max}
    \right]
  \label{eq:violation-rate}
\end{equation}


% -----------------------------------------------------------------------------
\subsection{Background: Soft Actor-Critic (SAC)}
\label{sec:temp-sac-background}

The CSAC-LB algorithm is built on Soft Actor-Critic
(SAC)~\cite{haarnoja2018sac}, an off-policy maximum-entropy actor-critic that
augments the standard RL objective with an entropy bonus:
\begin{equation}
  \pi^* = \arg\max_\pi \sum_{t=0}^{T}
    \mathbb{E}_{(s_t, a_t) \sim \rho_\pi}\!\left[
      \gamma^t \bigl(R(s_t, a_t) + \alpha\, H(\pi(\cdot|s_t))\bigr)
    \right]
  \label{eq:sac-objective}
\end{equation}
where $H(\pi(\cdot|s)) = -\mathbb{E}_{a \sim \pi}[\log \pi(a|s)]$ is the
policy entropy.  SAC maintains two Q-value networks (twin critics)
$Q_{\theta_1}, Q_{\theta_2}$ trained to minimise the soft Bellman residual:
\begin{align}
  y &= r + \gamma(1 - d_t)\!\left[
    \min_{j=1,2} Q_{\bar{\theta}_j}(s', a')
    - \alpha \log \pi_\phi(a'|s')
  \right], \quad a' \sim \pi_\phi(\cdot|s')
  \label{eq:sac-target} \\
  \mathcal{L}_Q(\theta_j) &= \mathbb{E}_{(s,a,r,s') \sim \mathcal{D}}\!\left[
    \bigl(Q_{\theta_j}(s,a) - y\bigr)^2
  \right]
  \label{eq:sac-critic-loss}
\end{align}
where $\mathcal{D}$ is the replay buffer, $\bar{\theta}$ are Polyak-averaged
target network parameters ($\bar{\theta} \leftarrow \tau\theta + (1-\tau)\bar{\theta}$),
and the $\min$ over twin critics prevents overestimation bias~\cite{fujimoto2018td3}.

The actor is a squashed Gaussian $\pi_\phi$: a neural network outputs
$(\mu_\phi(s), \log\sigma_\phi(s))$, actions are drawn via the
reparameterisation trick $a = \tanh(\mu_\phi(s) + \sigma_\phi(s) \odot \epsilon)$,
$\epsilon \sim \mathcal{N}(0, I)$, and the actor loss minimises:
\begin{equation}
  \mathcal{L}_\pi(\phi) = \mathbb{E}_{s \sim \mathcal{D}}\!\left[
    \alpha \log \pi_\phi(a|s)
    - \min_{j=1,2} Q_{\theta_j}(s, a)
  \right], \quad a \sim \pi_\phi(\cdot|s)
  \label{eq:sac-actor-loss}
\end{equation}

SAC is well-suited to building energy control for three reasons:
(i)~off-policy replay-buffer learning is sample-efficient when each training
episode spans a long seasonal horizon;
(ii)~the continuous stochastic policy handles the continuous $[0,1]^3$ cooling
action space naturally;
(iii)~entropy regularisation prevents collapse to a single cooling strategy,
enabling discovery of temporally varying control patterns that exploit price
and solar signals.


% -----------------------------------------------------------------------------
\subsection{CSAC-LB: Constrained SAC with Smoothed Log-Barrier}
\label{sec:temp-csac}

CSAC-LB~\cite{zhang2024csaclb} extends SAC with a constraint-handling mechanism
derived from interior-point methods.  We describe each component in detail,
since the algorithm is central to this chapter's results.

\subsubsection{The Constrained RL Problem and Lagrangian Limitations}

Standard SAC ignores the CMDP cost constraint.  The constrained optimisation
problem is:
\begin{equation}
  \max_\pi\; J_R(\pi) \quad \text{subject to} \quad J_C(\pi) \leq d
  \label{eq:cmdp-problem}
\end{equation}
The standard approach, SAC-Lag~\cite{ha2020learning}, converts this to a
min-max game via a Lagrange multiplier $\beta \geq 0$:
\begin{align}
  \max_\pi \min_{\beta \geq 0}\;
    L(\pi, \beta) &\doteq f(\pi) - \beta\, g(\pi)
  \label{eq:sac-lag-objective} \\
  f(\pi) &= \sum_{t} \mathbb{E}\!\left[\gamma^t
    \bigl(R(s_t,a_t) + \alpha H(\pi(\cdot|s_t))\bigr)\right]
  \label{eq:reward-objective} \\
  g(\pi) &= \sum_{t} \mathbb{E}\!\left[\gamma^t
    \bigl(C(a_t,s_t) - d\bigr)\right]
  \label{eq:cost-constraint}
\end{align}

Zhang et al.~\cite{zhang2024csaclb} identify three failure modes of SAC-Lag
that motivate the barrier approach:
\begin{enumerate}
  \item \textbf{Multiplier oscillation.}  The adversarial Lagrangian dynamics
    produce an unstable cycle: $\beta$ overshoots $\to$ overly conservative
    policy $\to$ cost drops $\to$ $\beta$ falls $\to$ policy becomes aggressive
    $\to$ cost rises.  This cycle prevents convergence near the constraint
    boundary and is confirmed empirically by our SAC-Lag runs (Section~\ref{sec:temp-results}),
    where the Lagrange multiplier grows monotonically from 3 to 27 without
    converging.
  \item \textbf{Dual-variable sensitivity.}  The stability of dual gradient
    descent depends critically on the relative learning rates of actor and
    multiplier.  An extra hyperparameter per scenario is required.
  \item \textbf{No boundary awareness.}  The Lagrangian multiplier applies a
    uniform penalty regardless of how close the policy is to the constraint
    boundary: it penalises equally a policy at $99\%$ of the limit and one at
    $1\%$.  This wastes capacity on excessive conservatism far from the boundary.
\end{enumerate}

\subsubsection{Linear Smoothed Log-Barrier Function}

CSAC-LB replaces the Lagrange multiplier with a differentiable barrier function
that is embedded directly in the actor loss.  The classical log-barrier
$-\frac{1}{\mu}\log(-x)$ for the constraint $x \leq 0$ is undefined when
$x > 0$ (infeasible), providing no recovery gradient when the policy has already
violated the constraint.  The \emph{linear smoothed} variant extends it to the
full real line:
\begin{equation}
  \bar{\psi}(x) =
  \begin{cases}
    -\dfrac{1}{\mu} \log(-x) & \text{if } x \leq -\dfrac{1}{\mu^2}
    \\[8pt]
    \mu\, x - \dfrac{1}{\mu}\log\!\left(\dfrac{1}{\mu^2}\right)
      + \dfrac{1}{\mu}
    & \text{otherwise}
  \end{cases}
  \label{eq:smoothed-barrier}
\end{equation}
where $\mu > 0$ controls the barrier steepness.

This construction has three critical properties:
\begin{enumerate}
  \item \textbf{$C^1$ differentiability everywhere.}  The two branches match
    in value and first derivative at $x = -1/\mu^2$, ensuring uninterrupted
    gradient flow through automatic differentiation.
  \item \textbf{Non-zero gradient in the infeasible region.}
    For $x > -1/\mu^2$ the gradient is the constant~$\mu$, providing a
    persistent \emph{recovery signal} that guides the policy back to feasibility.
    Clipped-penalty methods have zero gradient in this region, leaving the policy
    without directional information.
  \item \textbf{Proportional steepness near the boundary.}
    For $x \ll 0$ the gradient $\bar{\psi}'(x) = -1/(\mu x)$ increases as $x$
    approaches zero, creating an accelerating ``soft wall'' that slows the
    policy before the hard constraint.
\end{enumerate}

\paragraph{Applying the barrier to the CMDP cost constraint.}
To enforce $Q_C(s,a) \leq d$, the barrier is composed with a shifted ReLU:
\begin{equation}
  \tilde{\psi}^*(x) = \bar{\psi}\!\left(
    \operatorname{ReLU}(x - d) - \frac{1}{\mu^2}
  \right), \quad x = Q_C(s,a)
  \label{eq:barrier-shifted}
\end{equation}
The composition has three regimes:
\begin{itemize}
  \item $Q_C < d$ (constraint satisfied): $\operatorname{ReLU}(Q_C - d) = 0$,
    so the barrier input is $-1/\mu^2 < 0$; the penalty is near-zero.
  \item $Q_C \approx d$ (near-boundary): the penalty grows smoothly,
    creating a soft wall that prevents the policy from approaching the limit.
  \item $Q_C \gg d$ (violated): the penalty grows linearly with slope~$\mu$,
    providing a strong gradient to reduce cost.
\end{itemize}
The shift by $-1/\mu^2$ ensures the barrier activates slightly
\emph{before} the constraint boundary, providing an early-warning gradient.

\begin{figure}[htbp]
  \centering
  \includegraphics[width=0.45\textwidth]{figures/fig_barrier_function.pdf}
  \caption{Linear smoothed log-barrier $\bar\psi(x)$ for barrier steepnesses
    $\mu \in \{5, 15, 25, 50\}$.  Larger $\mu$ approximates the hard constraint
    more closely while maintaining $C^1$ differentiability.  The transition
    from logarithmic to linear behaviour occurs at $x = -1/\mu^2$, which
    moves toward zero as $\mu$ increases.  Adapted from
    Zhang et~al.~\cite{zhang2024csaclb}.}
  \label{fig:barrier-function}
\end{figure}

\subsubsection{Network Architecture: Six Specialised Networks}

CSAC-LB extends standard SAC (two reward critics, one actor) with two
additional cost critics and their corresponding target networks:
\begin{itemize}
  \item \textbf{Actor} $\pi_\phi$: Squashed Gaussian, outputs
    $(\mu_\phi, \log\sigma_\phi) \in \mathbb{R}^3 \times \mathbb{R}^3$.
  \item \textbf{Twin reward critics} $Q_{R,1}(\theta_{r_1}), Q_{R,2}(\theta_{r_2})$:
    Estimate discounted return.
  \item \textbf{Twin cost critics} $Q_{C,1}(\theta_{c_1}), Q_{C,2}(\theta_{c_2})$:
    Estimate discounted cost.
  \item \textbf{Target networks}: Slow-moving copies of all four critics,
    updated via Polyak averaging $\bar{\theta} \leftarrow \tau\theta +
    (1-\tau)\bar{\theta}$ with $\tau = 0.005$.
\end{itemize}
All networks use a two-layer MLP with $[256, 256]$ hidden units and ReLU
activations, taking the concatenated $[s, a] \in \mathbb{R}^{57}$ as input.

\subsubsection{Pessimistic Cost Critic: The Core Safety Mechanism}

The reward critics use $\min$ over the twin estimates (optimistic, preventing
reward overestimation):
\begin{equation}
  Q_R(s,a) = \min\bigl(Q_{R,1}(s,a),\; Q_{R,2}(s,a)\bigr)
  \label{eq:reward-critic-min}
\end{equation}
The cost critics use \emph{$\max$} (pessimistic, preventing cost
\emph{under}estimation):
\begin{equation}
  Q_C(s,a) = \max\bigl(Q_{C,1}(s,a),\; Q_{C,2}(s,a)\bigr)
  \label{eq:cost-critic-max}
\end{equation}
This asymmetry is the central safety mechanism.  In building control, most
replay-buffer transitions are constraint-satisfying ($C_t = 0$), creating a
systematic underestimation bias in the cost Q-function.  Taking the maximum of
two independently initialised critics counteracts this bias: the pessimistic
estimate is statistically more likely to flag constraint-risky actions as
unsafe, providing an implicit safety margin without requiring manual
tuning~\cite{zhang2024csaclb}.

The cost-critic targets differ from reward-critic targets in two ways:
\begin{equation}
  y_C = c_t + \gamma\,(1 - d_t)\,\max\!\left(
    Q_{\bar{C},1}(s', a'),\; Q_{\bar{C},2}(s', a')
  \right)
  \label{eq:cost-target}
\end{equation}
(i)~$\max$ is used for bootstrapping to propagate pessimism recursively; and
(ii)~the entropy term $-\alpha \log \pi$ is \emph{absent}: entropy
regularisation encourages exploration of high-entropy (uncertain) actions,
which is desirable for reward maximisation but undesirable for cost
estimation where prediction accuracy is the priority.

\subsubsection{Actor Loss with Barrier Penalty}

The complete actor loss combines standard SAC reward maximisation with the
barrier safety penalty, \emph{replacing the Lagrange multiplier entirely}:
\begin{equation}
  \mathcal{L}_\pi(\phi) = \mathbb{E}_{s \sim \mathcal{D}}\!\left[
    \underbrace{\alpha \log \pi_\phi(a|s)
      - \min(Q_{R,1}, Q_{R,2})}_{\text{entropy-regularised reward maximisation}}
    + \underbrace{\tilde{\psi}^*\!\bigl(
      \max(Q_{C,1}, Q_{C,2})\bigr)}_{\text{log-barrier safety penalty}}
  \right]
  \label{eq:csac-loss}
\end{equation}
The gradient of the barrier term with respect to the actor parameters flows as:
\begin{equation}
  \frac{\partial \tilde{\psi}^*}{\partial \phi} =
    \bar{\psi}'\!\bigl(\text{slack}\bigr) \cdot
    \frac{\partial\, \max(Q_{C,1}, Q_{C,2})}{\partial a} \cdot
    \frac{\partial a}{\partial \phi}
  \label{eq:barrier-gradient}
\end{equation}
where $\text{slack} = \operatorname{ReLU}(Q_C - d) - 1/\mu^2$.  When
$Q_C \ll d$, $\bar{\psi}' \approx 0$ and the actor optimises the reward
freely.  As $Q_C \to d$, $\bar{\psi}'$ grows, smoothly transferring the
optimisation objective from reward maximisation to constraint satisfaction.
This creates a \emph{proximity-aware} safety mechanism absent from SAC-Lag's
constant-multiplier approach.

\subsubsection{Training Procedure and Hyperparameters}

Algorithm~\ref{alg:csac-lb} presents the complete CSAC-LB training loop.

\begin{algorithm}[htbp]
\caption{CSAC-LB Training (adapted from Zhang et~al.~\cite{zhang2024csaclb})}
\label{alg:csac-lb}
\begin{algorithmic}[1]
\REQUIRE Reward critics $Q_{R,1}, Q_{R,2}$; cost critics $Q_{C,1}, Q_{C,2}$;
  actor $\pi_\phi$; target networks $\bar{Q}_{R}, \bar{Q}_{C}$;
  replay buffer $\mathcal{D}$; cost limit $d$; barrier factor $\mu$
\STATE Initialise all networks with Kaiming uniform weights
\STATE $\delta \gets 1/\mu^2$ \hfill $\triangleright$ \textit{barrier shift}
\FOR{each environment step $t = 1, 2, \ldots$}
  \IF{$t \leq t_{\text{warmup}}$}
    \STATE $a_t \gets \text{UniformRandom}(\mathcal{A})$
  \ELSE
    \STATE $a_t \sim \pi_\phi(\cdot | s_t)$
  \ENDIF
  \STATE Execute $a_t$, observe $r_t, c_t, s_{t+1}$; store in $\mathcal{D}$
  \IF{$t > t_{\text{warmup}}$}
    \STATE Sample mini-batch $\mathcal{B} \sim \mathcal{D}$; draw $a' \sim \pi_\phi(\cdot|s')$
    \STATE \textbf{// Reward critic update}
    \STATE $y_R \gets r + \gamma(1-d_t)[\min_j \bar{Q}_{R,j}(s',a') - \alpha \log \pi_\phi(a'|s')]$
    \STATE Minimise $\|Q_{R,j}(s,a) - y_R\|^2$ for $j=1,2$
    \STATE \textbf{// Cost critic update (pessimistic)}
    \STATE $y_C \gets c + \gamma(1-d_t)\,\max_j \bar{Q}_{C,j}(s',a')$
      \hfill \textit{max, no entropy}
    \STATE Minimise $\|Q_{C,j}(s,a) - y_C\|^2$ for $j=1,2$
    \IF{$t \bmod 2 = 0$} \hfill $\triangleright$ \textit{policy delay}
      \STATE \textbf{// Actor update with barrier penalty}
      \STATE $\tilde{a} \sim \pi_\phi(\cdot | s)$
      \STATE $Q_C \gets \max_j Q_{C,j}(s, \tilde{a})$
      \STATE $\text{slack} \gets \operatorname{ReLU}(Q_C - d) - \delta$
      \STATE $\mathcal{L}_\pi \gets \alpha \log \pi_\phi(\tilde{a}|s) - \min_j Q_{R,j}(s,\tilde{a}) + \bar{\psi}(\text{slack})$
      \STATE Minimise $\mathcal{L}_\pi$ w.r.t.\ $\phi$ (gradient clip: max norm $40.0$)
      \STATE Polyak-update all target networks
    \ENDIF
  \ENDIF
\ENDFOR
\end{algorithmic}
\end{algorithm}

Table~\ref{tab:csac-hyperparams} lists the CSAC-LB hyperparameters used in
all experiments.  The cost limit $d = 0.5$ and barrier steepness $\mu = 25$
were selected via grid search over $d \in \{0.1, 0.5, 1.0, 2.0, 5.0\}$ and
$\mu \in \{5, 10, 15, 25, 50\}$; the chosen combination achieved the best
violation-rate/energy-efficiency trade-off.  The entropy temperature is fixed
at $\alpha = 0.2$ (not auto-tuned) because automatic entropy tuning can
interfere with the barrier mechanism: when the barrier pushes the actor toward
conservative actions, the policy entropy decreases, which would trigger alpha
to increase---counteracting the constraint enforcement.

\begin{table}[htbp]
\centering
\caption{CSAC-LB hyperparameters for the temperature case study.}
\label{tab:csac-hyperparams}
\begin{tabular}{lcc}
\toprule
\textbf{Hyperparameter} & \textbf{Value} & \textbf{Source} \\
\midrule
\multicolumn{3}{l}{\textit{Network architecture}} \\
Hidden layers (all networks) & $[256, 256]$ & \cite{zhang2024csaclb} \\
Activation & ReLU & \cite{zhang2024csaclb} \\
Output squashing (actor) & $\tanh$ & \cite{haarnoja2018sac} \\
\midrule
\multicolumn{3}{l}{\textit{Optimisation}} \\
Learning rate & $10^{-3}$ & \cite{zhang2024csaclb} \\
Batch size & $256$ & \cite{zhang2024csaclb} \\
Discount factor $\gamma$ & $0.99$ & \cite{zhang2024csaclb} \\
Polyak $\tau$ & $0.005$ & \cite{zhang2024csaclb} \\
Policy delay & $2$ & \cite{zhang2024csaclb} \\
Gradient clip (max norm) & $40.0$ & \cite{zhang2024csaclb} \\
\midrule
\multicolumn{3}{l}{\textit{Replay buffer}} \\
Buffer capacity & $3 \times 10^6$ & Tuned \\
Random exploration steps & $100$ & Tuned \\
\midrule
\multicolumn{3}{l}{\textit{Constraint handling}} \\
Cost limit $d$ & $0.5$ & Tuned \\
Barrier steepness $\mu$ & $25.0$ & Tuned \\
\midrule
\multicolumn{3}{l}{\textit{Entropy and training}} \\
Entropy temperature $\alpha$ & $0.2$ (fixed) & \cite{zhang2024csaclb} \\
Training epochs & $60$ & Tuned \\
Steps per epoch & $2{,}208$ & Environment \\
\bottomrule
\end{tabular}
\end{table}


% -----------------------------------------------------------------------------
\subsection{Benchmark Algorithms}
\label{sec:temp-benchmark-algos}

Five safe RL algorithms are compared alongside the CSAC-LB primary method.
All algorithms share identical environment wrappers, observation normalisation,
and reward functions; the only differences are in the constraint-handling
mechanism and the policy update procedure.

\paragraph{CPO~\cite{achiam2017cpo}.}
Constrained Policy Optimisation solves a trust-region constrained subproblem at
each policy update: it linearises both the reward objective and the constraint
function around the current policy, then finds the maximal reward-improving
update within the constraint-feasible trust region.  The trust-region boundary
is a KL divergence constraint $\mathbb{E}[\text{KL}(\pi_{\text{old}} \| \pi)] \leq \delta_{\text{KL}}$,
ensuring stable policy improvement.  CPO's key strength is that the constraint
is enforced \emph{at every policy update step}, making it robust to
distributional shift.

\paragraph{FOCOPS~\cite{zhang2020focops}.}
First-Order Constrained Optimisation in Policy Space reformulates CPO's
constrained subproblem as a two-stage first-order procedure: (i)~unconstrained
optimisation of the reward objective to find an intermediate policy; (ii)~a
projection step that brings the intermediate policy back to the constraint
feasible set.  This avoids the expensive second-order computation of CPO's
trust region but introduces approximation error at the projection step.

\paragraph{CUP~\cite{yang2022cup}.}
Constrained Update Projection decomposes the policy gradient into reward-improving
and constraint-satisfying components, then combines them via a projection that
prioritises constraint satisfaction when the cost budget is exceeded.  Like
FOCOPS, CUP approximates CPO's per-step constraint guarantee using simpler
first-order machinery.

\paragraph{SAC-Lag~\cite{ha2020learning}.}
SAC with Lagrangian penalty adds a learned multiplier $\beta$ to the SAC
actor loss (Eq.~\ref{eq:sac-lag-objective}).  The multiplier is updated via
dual gradient descent.  Unlike CSAC-LB, there is no pessimistic cost critic
and no barrier function; the single multiplier must simultaneously detect
constraint violations and guide policy updates, creating the oscillation
failure mode described in Section~\ref{sec:temp-csac}.

\paragraph{RBC baseline.}
A deterministic bang-bang thermostat applies full cooling whenever any building
exceeds $T_{\text{bb}} = 24\,^\circ$C and zero cooling otherwise:
\begin{equation}
  a_t^{(i)} =
  \begin{cases}
    1.0 & T_t^{(i)} > T_{\text{bb}} \\
    0.0 & \text{otherwise}
  \end{cases}
  \label{eq:bangbang-rbc}
\end{equation}
This represents commodity HVAC thermostats and establishes a strong practical
baseline: it guarantees the temperature never rises above $T_{\text{bb}}$
between trigger events, but wastes energy through binary switching and produces
the highest ramping KPI (2.322) among all baselines.


% -----------------------------------------------------------------------------
\subsection{Evaluation Protocol and Training Configuration}
\label{sec:temp-eval-protocol}

\paragraph{Multi-seed training.}
Each algorithm is trained with \num{3}--\num{6} independent random seeds to
quantify variance arising from neural network initialisation and stochastic
policy rollouts.  CSAC-LB is trained on \num{6} seeds
($\{0, 1, 2, 7, 13, 42\}$); all on-policy algorithms (CPO, CUP, FOCOPS) and
SAC-Lag are trained on \num{3} seeds ($\{42, 0, 1\}$).  Training proceeds
sequentially to avoid OOM errors from concurrent replay buffer allocation
($3 \times 10^6$ transitions $\approx 2$--3\,GB per CSAC-LB instance).

\paragraph{Checkpoint selection.}
For each seed, two checkpoints are evaluated: the \emph{best} checkpoint
(the epoch with the lowest cumulative constraint cost during training) and the
\emph{final} checkpoint (epoch~60).  Best-checkpoint selection is based on
training-time cost monitoring, not held-out evaluation---it is therefore
available to practitioners without additional data.

\paragraph{Deterministic evaluation rollout.}
Each checkpoint is evaluated via a single deterministic rollout over the full
summer cooling-season horizon
(no exploration noise, greedy mean action for SAC-based algorithms): the
actor outputs $\hat{a}_t = \tanh(\mu_\phi(s_t))$ without stochastic sampling.
All metrics are computed over the 2{,}194 active steps.

\paragraph{Training configuration.}
All algorithms use $[256, 256]$ hidden layers, ReLU activations, discount
$\gamma = 0.99$, and the same \texttt{CityLearnTemp-CoolingOnly-Masked-Reward-v0}
environment registration.  On-policy algorithms (CPO, FOCOPS, CUP) use
2{,}208 steps per epoch and 60 total epochs; SAC-Lag uses a 200K-entry replay
buffer with a PID Lagrangian controller for stable multiplier updates.


% =============================================================================
\section{Evaluation Metrics}
\label{sec:temp-metrics}
% =============================================================================

We use the following key performance indicators (KPIs), consistent with the
CityLearn evaluation framework~\cite{nweye2023citylearn} and the STEMS
methodology~\cite{zhang2025stems}.

\paragraph{Comfort violation rate $V_{\text{rate}}$.}
The fraction of active time steps at which any building violates the comfort
band (Eq.~\ref{eq:violation-rate}).  This is the primary safety metric.

\paragraph{Discomfort proportion $D_{\text{prop}}$.}
The fraction of building-hour pairs outside the comfort band:
\begin{equation}
  D_{\text{prop}} = \frac{1}{T \cdot N}
    \sum_{t=1}^{T} \sum_{i=1}^{N}
    \mathbf{1}\!\left[T_{t}^{(i)} \notin [T_{\min}, T_{\max}]\right]
  \label{eq:discomfort-proportion}
\end{equation}

\paragraph{CityLearn KPI ratios.}
Electricity cost, carbon emissions, and consumption are reported as ratios
normalised relative to a no-storage baseline; values below 1.0 indicate
improvement.  Ramping, daily peak, and load factor are similarly normalised.


% =============================================================================
\section{Results and Discussion}
\label{sec:temp-results}
% =============================================================================

% -----------------------------------------------------------------------------
\subsection{Training Dynamics}
\label{sec:temp-results-training}

Figure~\ref{fig:training-curves} shows training trajectories across all
benchmark algorithms.  Three distinct convergence patterns emerge:

\textbf{Effective constraint learners (CSAC-LB, CPO):}  CSAC-LB drives the
episode cost from $\sim$280 to near-zero within the first 15 epochs, while the
episode return rises from $-4{,}000$ to $+2{,}000$.  The barrier penalty peaks
around epoch~10 as the cost critic detects constraint-approaching actions, then
decreases as the policy retreats from the constraint boundary.  CPO follows a
similar cost reduction curve, reaching near-zero cost by epoch~40.

\textbf{Moderate learners (SAC-Lag):}  SAC-Lag's Lagrange multiplier grows
monotonically from 3 to 27 over 60 epochs without converging, confirming the
oscillation failure mode described in Section~\ref{sec:temp-csac}.  The cost
critic value rises continuously to 8.5, never stabilising.

\textbf{Failed constraint learners (FOCOPS, CUP):}  Both algorithms plateau
at episode costs of 200--250 per epoch, two to three orders of magnitude above
the CSAC-LB/CPO regime.  Their mean cooling actions collapse to 0.102 (FOCOPS)
and 0.142 (CUP), indicating the policies converge to near-zero cooling, which
paradoxically reduces energy KPIs while causing 85--95\% comfort violations.

\begin{figure}[htbp]
  \centering
  \includegraphics[width=\textwidth]{figures/fig_training_diagnostics.pdf}
  \caption{Training diagnostics over 60 epochs for all five benchmark
    algorithms.  (a)~Episode return: CSAC-LB and CPO show clear upward
    learning curves.  (b)~Policy entropy: on-policy algorithms gradually
    decrease from 1.4 to 1.0--1.2 nats (healthy exploration-to-exploitation
    transition); CSAC-LB maintains fixed $\alpha=0.2$.  (c)~Episode cost
    (primary safety diagnostic): CSAC-LB and CPO drive cost toward zero;
    FOCOPS and CUP plateau at 200--250.  (d)~Value loss: CSAC-LB achieves
    the lowest final critic loss.  (e)~Policy loss: CSAC-LB's entropy-regularised
    policy loss decreases from +18 to -47, characteristic of SAC-style
    improvement.}
  \label{fig:training-curves}
\end{figure}


% -----------------------------------------------------------------------------
\subsection{Cooling-Season Evaluation: Multi-Algorithm Comparison}
\label{sec:temp-results-fullyr}

Table~\ref{tab:temp-full-benchmark} presents the cooling-season evaluation results
at the \emph{best checkpoint} for all algorithms, sorted by violation rate.
Two deterministic baselines---the bang-bang RBC and the Zero-cooling
controller---establish the practical performance range.

\begin{table}[htbp]
\centering
\caption{Cooling-season evaluation: best checkpoint results sorted by violation
  rate.  KPI ratios are normalised to the no-storage baseline; values $<1.0$
  indicate improvement.  Algorithms achieving $V_{\text{rate}} < 0.20$ are
  highlighted.  All values are from deterministic evaluation rollouts.}
\label{tab:temp-full-benchmark}
\begin{tabular}{lcccccc}
\toprule
\textbf{Algorithm}
  & \textbf{$V_{\text{rate}}$}
  & \textbf{$D_{\text{prop}}$}
  & \textbf{Cost KPI}
  & \textbf{Emis.\ KPI}
  & \textbf{Elec.\ KPI}
  & \textbf{Ramp KPI} \\
\midrule
\textbf{CSAC-LB (s1, best)} & \textbf{0.056} & 0.153 & 0.938 & 0.969 & 0.968 & 1.067 \\
CSAC-LB (s0, best)          & 0.068          & 0.461 & 1.109 & 1.156 & 1.153 & 1.015 \\
CPO (best)                  & 0.132          & 0.134 & 1.083 & 1.098 & 1.093 & 0.988 \\
CSAC-LB (s42, best)         & 0.145          & 0.079 & 0.857 & 0.878 & 0.875 & 1.065 \\
CSAC-LB (s2, best)          & 0.184          & 0.113 & 0.940 & 0.960 & 0.959 & 1.075 \\
\midrule
SAC-Lag (best)              & 0.327          & 0.065 & 0.779 & 0.800 & 0.798 & 1.403 \\
\midrule
RBC                         & 0.516          & 0.376 & 1.010 & 1.043 & 1.041 & 2.322 \\
PPO (best)                  & 0.428          & 0.060 & 0.759 & 0.780 & 0.779 & 1.282 \\
\midrule
CUP (best)                  & 0.856          & 0.610 & 0.843 & 0.817 & 0.816 & 1.125 \\
FOCOPS (best)               & 0.953          & 0.847 & 0.667 & 0.692 & 0.688 & 1.189 \\
Zero                        & 0.995          & 0.962 & 0.527 & 0.562 & 0.559 & 0.905 \\
\bottomrule
\end{tabular}
\end{table}

\paragraph{CSAC-LB achieves the lowest violation rates.}
CSAC-LB seed~1 attains $V_{\text{rate}} = 0.056$---an 89.1\% reduction
relative to the RBC baseline ($0.516$).  Four of the six CSAC-LB seeds achieve
$V_{\text{rate}} < 0.185$ at best checkpoint, all below the Cooling RBC.
CPO achieves $0.132$, also substantially below the RBC.

\paragraph{The safety--efficiency trade-off.}
CSAC-LB seed~42 lies near the empirical Pareto frontier: $V_{\text{rate}} =
0.145$ with a cost KPI of $0.857$---a 14.3\% energy saving below the
no-storage baseline while maintaining comparatively low violation rates.
FOCOPS achieves the
lowest cost KPI ($0.667$) but with 95.3\% violation rate, because it
consistently under-cools.  No algorithm simultaneously achieves the lowest cost
and the lowest violation rate, confirming the theoretical prediction that
safety and efficiency are competing objectives.

\paragraph{Reward decomposition.}
Table~\ref{tab:reward-decomposition} decomposes cumulative reward into its four
STEMS components.  The comfort penalty dominates the total: FOCOPS receives
$R^{\text{comfort}} = -4{,}100.8$ despite competitive stability and renewable
scores, yielding a net negative total reward.  CSAC-LB seeds achieve
$R^{\text{comfort}}$ values between $-17.0$ and $-80.7$, confirming effective
temperature regulation.

\begin{table}[htbp]
\centering
\caption{Cumulative STEMS reward components at best checkpoint.
  Higher (less negative) is better in all cases.}
\label{tab:reward-decomposition}
\begin{tabular}{lrrrrc}
\toprule
\textbf{Algorithm}
  & $R^{\text{econ}}$
  & $R^{\text{stab}}$
  & $R^{\text{renew}}$
  & $R^{\text{comfort}}$
  & \textbf{Total} \\
\midrule
CSAC-LB (s1)  & $-251.9$ & $1{,}591.1$ & $946.0$   & $-17.0$    & $2{,}268.2$ \\
CSAC-LB (s42) & $-223.6$ & $2{,}065.6$ & $1{,}011.8$ & $-80.7$  & $2{,}773.2$ \\
CPO           & $-260.2$ & $1{,}490.1$ & $756.9$   & $-55.3$    & $1{,}931.4$ \\
SAC-Lag       & $-178.5$ & $2{,}043.2$ & $1{,}120.8$ & $-234.3$ & $2{,}751.3$ \\
CUP           & $-196.2$ & $2{,}502.1$ & $1{,}020.1$ & $-3{,}163.8$ & $162.1$ \\
FOCOPS        & $-165.1$ & $2{,}730.2$ & $1{,}162.2$ & $-4{,}100.8$ & $-373.4$ \\
RBC           & $-229.1$ & $1{,}703.6$ & $1{,}136.6$ & $-3{,}059.7$ & $-448.7$ \\
Zero          & $-93.0$  & $3{,}403.5$ & $1{,}416.5$ & $-50{,}158.6$ & $-45{,}431.6$ \\
\bottomrule
\end{tabular}
\end{table}


% -----------------------------------------------------------------------------
\subsection{Multi-Seed Statistical Analysis}
\label{sec:temp-results-multiseed}

To rigorously quantify algorithmic robustness, each algorithm was trained on
\num{3}--\num{6} random seeds and evaluated with the best-checkpoint protocol.
Table~\ref{tab:multiseed-summary} reports mean $\pm$ standard deviation for
the three primary KPIs.

\begin{table}[htbp]
\centering
\caption{Multi-seed evaluation summary: mean $\pm$ standard deviation at
  best checkpoint.  CSAC-LB is evaluated across $n=6$ seeds; all other
  algorithms across $n=3$ seeds.  Lower violation rate and lower cost are
  better; higher reward is better.  \textbf{Bold} denotes the best mean
  violation rate.}
\label{tab:multiseed-summary}
\begin{tabular}{lcrrrr}
\toprule
\textbf{Algorithm} & $n$ &
  \textbf{Reward $\mu \pm \sigma$} &
  \textbf{Cost $\mu \pm \sigma$} &
  \textbf{Viol.\ Rate $\mu \pm \sigma$} \\
\midrule
\textbf{CSAC-LB} & 6 &
  $2{,}348.4 \pm 321.7$ &
  $3.94 \pm 1.69$ &
  $\mathbf{0.113 \pm 0.050}$ \\
SAC-Lag & 3 &
  $2{,}148.4 \pm 6.6$ &
  $25.39 \pm 1.02$ &
  $0.205 \pm 0.008$ \\
CPO & 3 &
  $1{,}753.3 \pm 171.0$ &
  $16.91 \pm 16.81$ &
  $0.316 \pm 0.239$ \\
FOCOPS & 3 &
  $1{,}979.4 \pm 207.2$ &
  $225.58 \pm 31.47$ &
  $0.931 \pm 0.026$ \\
CUP & 3 &
  $1{,}828.6 \pm 334.0$ &
  $216.54 \pm 66.47$ &
  $0.919 \pm 0.032$ \\
\midrule
\multicolumn{5}{l}{\textit{Baselines (deterministic, single run)}} \\
Cooling RBC & -- & $-448.7$ & -- & $0.516$ \\
Zero & -- & $-45{,}431.6$ & -- & $0.995$ \\
\bottomrule
\end{tabular}
\end{table}

\paragraph{CSAC-LB statistical robustness.}
Across six seeds, CSAC-LB achieves $V_{\text{rate}} = 0.113 \pm 0.050$ with
a coefficient of variation $\text{CV} = 0.050/0.113 = 0.44$.  The reward
standard deviation ($\sigma = 321.7$, $\text{CV} = 0.14$) is moderate but
acceptable; all six seeds converge to violation rates substantially below the
Cooling RBC (0.516).  The low cost standard deviation ($\sigma = 1.69$,
$\text{CV} = 0.43$) confirms that the barrier mechanism reliably enforces
constraints across diverse random initialisations.

\paragraph{Per-seed breakdown.}
Table~\ref{tab:csaclb-perseed} details the six individual CSAC-LB seeds.
Seed~1 achieves the best-checkpoint violation rate of $0.056$ (epoch~29),
while seed~2 requires the most training epochs (epoch~53) due to a slower
cost reduction trajectory.  Best-epoch selection varies between epoch~24
(seed~0) and epoch~59 (seed~42), illustrating that cost-based checkpoint
monitoring is essential for achieving peak safety performance.

\begin{table}[htbp]
\centering
\caption{CSAC-LB per-seed results.  Best checkpoint is selected by minimum
  training cost.  Total reward is from the deterministic cooling-season rollout.}
\label{tab:csaclb-perseed}
\begin{tabular}{crrccc}
\toprule
\textbf{Seed} &
  \textbf{Best Ep.} &
  \textbf{Viol.\ Rate} &
  \textbf{Total Cost} &
  \textbf{Total Reward} &
  \textbf{Discomfort} \\
\midrule
0  & 24 & 0.068 & 2.046 & 1{,}811.4 & 0.461 \\
1  & 29 & \textbf{0.056} & 1.976 & 2{,}268.2 & 0.153 \\
2  & 53 & 0.184 & 6.024 & 2{,}548.6 & 0.113 \\
7  & -- & 0.089 & 3.725 & 2{,}377.2 & 0.215 \\
13 & -- & 0.136 & 5.445 & 2{,}312.1 & 0.232 \\
42 & 59 & 0.145 & 4.444 & 2{,}773.2 & 0.079 \\
\midrule
\textbf{Mean} & -- & $0.113$ & $3.94$ & $2{,}348.4$ & -- \\
\textbf{Std} & -- & $0.050$ & $1.69$ & $321.7$ & -- \\
\bottomrule
\end{tabular}
\end{table}

\paragraph{SAC-Lag: high reward, moderate safety.}
SAC-Lag achieves the tightest reward variance across seeds ($\sigma = 6.6$,
$\text{CV} = 0.003$), confirming that off-policy SAC's experience replay
produces reproducible reward convergence.  However, its violation rate
($0.205 \pm 0.008$) is nearly double that of CSAC-LB, confirming that the
Lagrangian multiplier alone is insufficient for tight constraint satisfaction.

\paragraph{CPO: high variance under multi-seed evaluation.}
CPO achieves a mean violation rate of $0.316 \pm 0.239$, with seed~42
achieving $0.088$ while seeds~0 and~1 converge to $0.294$ and $0.565$
respectively.  This large variance ($\text{CV} = 0.76$) suggests that CPO's
trust-region mechanism is sensitive to random initialisation and that the
seed~42 single-run result reported in prior single-seed benchmarks may be
optimistic.  The high variance in cost ($16.91 \pm 16.81$) further confirms
this instability.

\paragraph{FOCOPS and CUP: consistent failure across seeds.}
Both algorithms show low violation-rate variance ($\text{CV} = 0.028$ for
FOCOPS, $0.035$ for CUP), but this consistency reflects \emph{reliably poor}
performance: every seed converges to near-zero cooling policies with violation
rates of 85--95\%.  The first-order and projection-based constraint
approximations of these methods fail to correctly identify the
constraint-satisfying policy direction in the presence of thermal inertia and
multi-building coupling.

\begin{figure}[htbp]
  \centering
  \includegraphics[width=0.80\textwidth]{figures/fig_multiseed_robustness.pdf}
  \caption{Multi-seed violation rate comparison across all five benchmark
    algorithms.  Each point represents one seed's best-checkpoint violation
    rate.  Horizontal bars indicate the mean.  CSAC-LB ($n=6$) achieves the
    lowest mean violation rate; CPO ($n=3$) shows the highest inter-seed
    variance; FOCOPS and CUP ($n=3$) cluster near 90\% violation despite
    small variance.  The dashed line marks the Cooling RBC baseline (51.6\%).}
  \label{fig:multiseed-robustness}
\end{figure}


% -----------------------------------------------------------------------------
\subsection{Behavioral Cluster Analysis}
\label{sec:temp-results-clusters}

The multi-seed results confirm three behavioural clusters, visible in both the
training dynamics and the evaluation KPIs.

\paragraph{Cluster~1: Effective constraint learners (CSAC-LB, CPO).}
These algorithms achieve mean violation rates below 20\% and demonstrate
genuine constraint-aware policy optimization.
\begin{itemize}
  \item \textbf{CSAC-LB} uses a conservative lower-bound safety critic that
    provides a pessimistic estimate of constraint cost, maintaining a safety
    margin that accounts for estimation uncertainty.  The off-policy replay
    buffer enables efficient learning from diverse state-action-cost tuples.
  \item \textbf{CPO} uses a trust-region approach that directly constrains
    expected cost at each update, providing a structural safety guarantee
    that is robust to moderate distributional shift.  Its main limitation is
    high seed variance (Table~\ref{tab:multiseed-summary}).
\end{itemize}

\paragraph{Cluster~2: Moderate constraint learner (SAC-Lag).}
SAC-Lag reduces violations relative to the RBC but at $0.205$ violation rate
cannot achieve the tight constraint satisfaction of CSAC-LB.  Its reward
consistency ($\sigma = 6.6$) demonstrates SAC's off-policy stability, but the
Lagrangian mechanism introduces the boundary-agnostic penalty failure mode
that prevents convergence near the constraint boundary.

\paragraph{Cluster~3: Failed constraint learners (FOCOPS, CUP).}
Both algorithms converge to near-zero cooling policies ($\bar{a} = 0.10$--$0.14$)
with violation rates exceeding 85\%.  The failure mechanism appears to be a
mismatch between the algorithms' first-order constraint approximations and the
CityLearn environment's nonlinear thermal dynamics:
\begin{itemize}
  \item \textbf{Thermal inertia:} building temperature responds to cooling
    actions with a one-step delay; first-order methods assume a more immediate
    cost-action relationship.
  \item \textbf{Multi-building coupling:} the district-level cost depends on
    the joint action across all three buildings; FOCOPS's two-stage projection
    may fail to find the correct projection direction in this coupled space.
\end{itemize}


% -----------------------------------------------------------------------------
\subsection{Heat-Wave Stress Test}
\label{sec:temp-results-heatwave}

To evaluate robustness under distributional shift, we use the heat-wave
evaluation artifact provided with this case study: a 72-hour extreme-temperature
scenario derived from the hottest contiguous summer period in the dataset.
Following the evaluation scripts, cooling demand is amplified by $2.5\times$
and solar generation by $1.2\times$ to simulate extreme conditions.  Only the
best checkpoints of each algorithm are evaluated.

\begin{table}[htbp]
\centering
\caption{72-hour heat-wave stress test: best checkpoint results, sorted by
  violation rate.  The cooling demand multiplier is $2.5\times$ the nominal
  value.  Bold indicates the best value per metric.}
\label{tab:heatwave}
\begin{tabular}{lccccc}
\toprule
\textbf{Algorithm}
  & \textbf{$V_{\text{rate}}$}
  & \textbf{$D_{\text{prop}}$}
  & \textbf{Reward}
  & \textbf{Total Cost}
  & \textbf{Mean Action} \\
\midrule
\textbf{CPO}        & \textbf{0.026} & 0.180 & 103.2  & 0.010 & 0.333 \\
CSAC-LB (s1)        & 0.066          & 0.330 & 110.6  & 0.051 & 0.297 \\
CSAC-LB (s1, final) & 0.066          & 0.170 & 120.4  & 0.035 & 0.256 \\
CPO (final)         & 0.105          & 0.146 & 106.0  & 0.078 & 0.316 \\
\midrule
Cooling RBC         & 0.540          & 0.481 & 81.0   & 5.389 & 0.348 \\
\midrule
CUP (best)          & 0.763          & 0.500 & 49.8   & 10.297 & 0.231 \\
FOCOPS (best)       & 0.803          & 0.490 & 63.7   & 11.123 & 0.173 \\
SAC-Lag (best)      & 0.947          & 0.413 & 56.6   & 9.009  & 0.243 \\
Zero                & 1.000          & 0.976 & $-467$ & 90.567 & 0.000 \\
\bottomrule
\end{tabular}
\end{table}

\paragraph{CPO performs best under the stress-test protocol.}
CPO attains $V_{\text{rate}} = 0.026$ under $2.5\times$ cooling demand,
indicating comparatively high robustness to distributional shift.  This reverses
the cooling-season ranking (where CSAC-LB
leads) and is explained by CPO's trust-region mechanism: the per-update
constraint guarantee is \emph{structural} rather than learned, so it cannot
become miscalibrated when the data distribution shifts to an extreme heat wave.
CPO's higher mean cooling action ($0.333$ vs.\ CSAC-LB's $0.297$) suggests it
learns a more aggressive cooling strategy that provides greater headroom under
extreme conditions.

\paragraph{CSAC-LB remains competitive under stress.}
CSAC-LB (seed~1) achieves $V_{\text{rate}} = 0.066$ under stress---only
marginally above its cooling-season performance ($0.056$)---demonstrating that the
pessimistic cost critic's safety margin is robust to extreme conditions.

\paragraph{SAC-Lag degrades catastrophically.}
SAC-Lag's violation rate rises from $0.205$ (cooling season) to $0.947$ (stress
test), confirming the theoretical prediction that Lagrangian multipliers can
become miscalibrated under distributional shift: the multiplier learned during
normal operation provides insufficient correction when the environment's cost
dynamics change abruptly.


% -----------------------------------------------------------------------------
\subsection{CityLearn KPI Analysis}
\label{sec:temp-results-kpi}

Table~\ref{tab:citylearn-kpi} compares the full CityLearn KPI suite for the
top-performing algorithms against the Cooling RBC.

\begin{table}[htbp]
\centering
\caption{Full CityLearn KPI comparison at best checkpoint.
  Values $<1.0$ indicate improvement over the no-storage baseline.
  CSAC-LB achieves large ramping improvement (53\%--54\%) across all
  seeds, while CPO achieves the best all-time peak KPI (0.965).}
\label{tab:citylearn-kpi}
\begin{tabular}{lccccccc}
\toprule
\textbf{Metric} &
  \textbf{CSAC-LB (s1)} &
  \textbf{CSAC-LB (s42)} &
  \textbf{CPO} &
  \textbf{SAC-Lag} &
  \textbf{RBC} \\
\midrule
Cost Total            & 0.938 & 0.857 & 1.083 & 0.779 & 1.010 \\
Carbon Emissions      & 0.969 & 0.878 & 1.098 & 0.800 & 1.043 \\
Electricity Cons.     & 0.968 & 0.875 & 1.093 & 0.798 & 1.041 \\
Ramping               & 1.067 & 1.065 & 0.988 & 1.403 & 2.322 \\
Daily Peak            & 1.071 & 1.041 & 1.052 & 0.889 & 1.131 \\
All-Time Peak         & 1.008 & 1.006 & 0.965 & --    & 1.059 \\
Discomfort Prop.      & 0.153 & 0.079 & 0.134 & 0.065 & 0.376 \\
\textbf{Viol.\ Rate}  & \textbf{0.056} & \textbf{0.145} & 0.132 & 0.327 & 0.516 \\
\bottomrule
\end{tabular}
\end{table}

The most striking result across all algorithms is the ramping improvement.
All effective safe RL agents dramatically outperform the RBC on ramping
(1.067 vs.\ 2.322 for CSAC-LB~s1, 0.988 vs.\ 2.322 for CPO), because the
bang-bang controller's binary switching creates large load spikes.  RL agents
learn smooth continuous cooling profiles that inherently avoid ramping.

CSAC-LB seed~42 achieves the best cost KPI ($0.857$) and emissions KPI
($0.878$), representing 14.3\% and 12.2\% improvements respectively over
the no-storage baseline---simultaneously with a $0.145$ violation rate.  This
confirms CSAC-LB's ability to navigate the safety--efficiency trade-off near
the Pareto frontier.


% -----------------------------------------------------------------------------
\subsection{Discussion}
\label{sec:temp-discussion}

\paragraph{Why CSAC-LB outperforms on average.}
CSAC-LB's superior mean violation rate arises from three synergistic
architectural advantages:
\begin{enumerate}
  \item \textbf{Pessimistic cost critic:} The $\max$ aggregation over twin
    cost critics provides a principled safety margin that accounts for
    function approximation uncertainty, particularly important in building
    control where constraint-violating transitions are rare in the replay
    buffer.
  \item \textbf{Off-policy sample efficiency:} The 3M-entry replay buffer
    enables data reuse across all historical transitions, allowing the cost
    critic to learn from a diverse set of constraint-approaching states even
    when the current policy rarely visits them.
  \item \textbf{Barrier proximity awareness:} The log-barrier gradient grows
    as $Q_C \to d$, creating a soft wall that steers the policy away from
    the constraint boundary without requiring a learned dual variable.
\end{enumerate}

\paragraph{Why CPO excels under distributional shift.}
CPO's structural trust-region constraint provides per-update guarantees that
cannot become miscalibrated when the environment dynamics change.  Under the
$2.5\times$ heat-wave stress, this structural property translates to a
$0.026$ violation rate---4.3$\times$ better than CSAC-LB.  For
safety-critical deployments where worst-case performance under extreme
conditions is the primary concern, CPO is the preferred choice.

\paragraph{Implications for practical deployment.}
The multi-seed analysis reveals that single-seed benchmarking can be misleading:
CPO's seed~42 achieves $0.088$ violation rate (comparable to CSAC-LB), but its
mean across three seeds is $0.316$.  A practical recommendation for building
energy management system (BEMS) deployors is to train each algorithm on at
least three seeds and use best-checkpoint selection based on training-time cost
monitoring; single-run results may significantly overstate reliability.


% =============================================================================
\section{Summary}
\label{sec:temp-summary}
% =============================================================================

This chapter presented a comprehensive multi-algorithm, multi-seed benchmark
for safe RL in HVAC temperature control.  The key findings are:

\begin{enumerate}
  \item \textbf{CSAC-LB achieves the best safety--efficiency trade-off.}
    Across six seeds, mean violation rate $0.113 \pm 0.050$ (best checkpoint),
    with CSAC-LB~s1 achieving $0.056$---an 89.1\% reduction relative to the
    Cooling RBC.  Cost and emissions KPIs of 0.857--0.938 demonstrate
    simultaneous energy efficiency.

  \item \textbf{CPO is the most robust algorithm under distributional shift.}
    2.6\% violation rate under the 72-hour, $2.5\times$ heat-wave stress test,
    outperforming CSAC-LB ($6.6\%$) and all other algorithms.  CPO's structural
    trust-region mechanism cannot be miscalibrated by distributional shift.

  \item \textbf{SAC-Lag provides moderate safety but lacks robustness.}
    Mean violation rate $0.205 \pm 0.008$ (tight, consistent), but degrades
    to $0.947$ under the heat-wave stress test.  The Lagrangian multiplier
    oscillation failure mode is consistent with the monotonically growing
    multiplier trajectory.

  \item \textbf{FOCOPS and CUP fail to learn effective constraint
    satisfaction.}  Consistent 85--95\% violation rates across all seeds,
    arising from a fundamental mismatch between their first-order constraint
    approximations and the environment's nonlinear thermal dynamics and
    multi-building coupling.

  \item \textbf{Multi-seed evaluation is essential.}  Single-seed benchmarks
    can overstate safety performance (CPO seed~42: $0.088$, mean: $0.316$).
    Best-checkpoint selection based on training-time cost monitoring is
    critical for achieving peak safety, since the best epoch varies widely
    across seeds (epoch~24 to epoch~59 for CSAC-LB).
\end{enumerate}

These findings have direct implications for BEMS deployment.  CSAC-LB is
recommended as the primary algorithm when both comfort safety and energy
efficiency are required.  CPO should be preferred when robustness to extreme
weather events is the primary operational concern.  Lagrangian methods (SAC-Lag,
PPO-Lag) provide a useful middle ground when computational constraints limit
replay-buffer memory, but should be validated against extreme-condition stress
tests before deployment.
